\documentclass[12pt]{article}
\usepackage[utf8]{inputenc}
\usepackage{geometry}
\geometry{letterpaper, margin=1in}
\usepackage{hyperref}

\title{\textbf{User Manual \& README}\\ \large EagleNav Project}
\author{Moh's Eagles}
\date{May 15, 2026}

\begin{document}

\maketitle

\section{Jira Sprint Board}
\textbf{Link:} \url{https://calstatela-cs3338-eaglenav.atlassian.net/jira/software/projects/KAN/boards/2}

\section{Formal Objective Breakdown}
EagleNav is an accessible, augmented reality campus navigation platform. The primary objectives are:
\begin{enumerate}
    \item Provide highly accurate, closest-entrance pedestrian routing for the CSULA campus.
    \item Deliver multimodal guidance (audio, haptic, visual) to support students with visual or physical impairments.
    \item Integrate real-time campus events and spatial awareness features.
    \item Operate with strict data privacy, keeping computations and location data on-device where possible.
\end{enumerate}

\section{Why This Software Matters}
Standard navigation apps like Google Maps route users to the "centroid" or middle of a building, often leaving visually impaired students stranded in parking lots or facing inaccessible staircases. EagleNav matters because it utilizes high-precision local data to route users specifically to ADA-accessible doors, removing physical barriers to education and fostering campus independence.

\section{Installation \& Access Instructions}
\subsection{Prerequisites}
\begin{itemize}
    \item Docker and Docker Compose installed on your host machine.
    \item Flutter SDK (v3.0+) and Dart installed for the mobile client.
    \item Android Studio (for Android emulation) or Xcode (for iOS emulation).
\end{itemize}

\subsection{Backend Setup (Docker)}
\begin{enumerate}
    \item Clone the GitHub repository to your local machine.
    \item Open a terminal or Visual Studio developer command prompt in the root directory.
    \item Run the command: \texttt{docker-compose up -d}
    \item The Valhalla routing engine will initialize on port 8002, and the simulated Node.js API will initialize on port 3000.
\end{enumerate}

\subsection{Frontend Setup (Mobile App)}
\begin{enumerate}
    \item Navigate to the \texttt{/mobile} directory in the repository.
    \item Run \texttt{flutter pub get} to install all required dependencies (e.g., flutter\_map, tflite).
    \item Connect a physical mobile device or launch an emulator.
    \item Run \texttt{flutter run} to build and launch the application.
\end{enumerate}

\end{document}
